\begin{frame}{손으로 한 번: 베이스라인의 F1}
  8.2절 사기 탐지 데이터(99.8:0.2)에서 다수 클래스(정상)만
  찍으면:
  \begin{itemize}
    \item 부정(fraud)에 대한 정확도는 $P=1.0$ 이지만 \textbf{재현율
      $R=0$} --- 부정 건을 하나도 못 잡음.
    \item \[ F_1 = \frac{2PR}{P+R} = 0 \]
  \end{itemize}
  반대로 부정만 전부 맞히려는 극단(전부를 부정으로 찍기)에서는:
  \begin{itemize}
    \item 재현율은 $R=1.0$ 이지만 정확도는 $0.2\%$ 수준으로 희석.
    \item \[ F_1 \approx 0.004 \]
  \end{itemize}
  \vskip0.4em
  \begin{block}{교훈}
    F1은 정확도보다 이런 \textbf{극단적 모델을 더 가혹하게
    (= 정확하게)} 다뤄준다 --- 바로 이것이 지표를 고르는
    이유다.
  \end{block}
\end{frame}
